\section{From cumulative bound time to fractional occupancy}
\label{sec:occupancy}

\subsection{Classical two-state reference}

For independent receptors exposed to a spatially uniform bulk concentration
$C$, classical mass-action kinetics gives
\begin{equation}
 \frac{d\theta}{dt}=\kOn C(1-\theta)-\kOff\theta.
 \label{eq:langmuir-kinetics}
\end{equation}
At steady state,
\begin{equation}
 \theta_{\mathrm{classical}}(C)
 =\frac{\kOn C}{\kOn C+\kOff}
 =\frac{C}{C+K_D},
 \qquad K_D=\frac{\kOff}{\kOn}.
 \label{eq:langmuir-isotherm}
\end{equation}
This is the noncooperative Langmuir reference for a one-to-one binding site
\citep{langmuir1918}.  It assumes that dissociation ends the occupied episode;
the same ligand's subsequent return is not retained as memory in the
two-state rate equation.

\subsection{Renewal--reward closure with rebinding}

Let
\begin{equation}
 \tauB(\phi)=\E[T_B\mid\text{background occupancy }\phi]
 \label{eq:effective-bound-time}
\end{equation}
be the mean cumulative time for which a tagged ligand is bound before it
escapes the near-surface region.  It includes the initial bound interval and
all subsequent bound intervals, but excludes intervening free-diffusion time.
Thus, $\tauB$ is the total bound-state ``reward'' generated by one successful
arrival into the retained trajectory.

For $N_R$ statistically equivalent receptor sites with ensemble occupancy
$\theta$, let $\lambda_{\mathrm{in}}$ denote the total rate at which ligands
associate with vacant receptors and initiate retained trajectories.  Under
the mean-field association model,
\begin{equation}
 \lambda_{\mathrm{in}}
 =N_R\kOn C(1-\theta).
 \label{eq:arrival-rate}
\end{equation}
The expected number of bound receptors at steady state is the trajectory
initiation rate multiplied by the mean cumulative bound time generated by
each trajectory,
\begin{equation*}
 \left\langle N_B\right\rangle
 =\lambda_{\mathrm{in}}\tauB(\theta).
\end{equation*}
By the definition of fractional occupancy, the same expected number is
\begin{equation*}
 \left\langle N_B\right\rangle=N_R\theta.
\end{equation*}
Combining these two identities with Eq.~\ref{eq:arrival-rate} gives
\begin{equation*}
 N_R\theta=N_R\kOn C(1-\theta)\tauB(\theta).
\end{equation*}
Cancelling $N_R$ gives the occupancy balance
\begin{equation}
 \theta=\kOn C(1-\theta)\tauB(\theta).
 \label{eq:renewal-balance}
\end{equation}
For an occupancy-independent effective lifetime, Eq.~\ref{eq:renewal-balance}
can be written
\begin{equation}
 \theta(C)=\frac{\kOn C\tauB}{1+\kOn C\tauB}.
 \label{eq:effective-occupancy}
\end{equation}
This is equivalently a two-state kinetic expression with operational removal
rate $k_{\mathrm{off,eff}}=1/\tauB$.  In the continuous-time no-rebinding
limit, $\tauB=1/\kOff$, and Eq.~\ref{eq:effective-occupancy} reduces to the
classical result in Eq.~\ref{eq:langmuir-isotherm}.  In the implemented
discrete process, the corresponding no-rebinding mean is given by
Eq.~\ref{eq:discrete-dwell-mean} and differs from $1/\kOff$ by a relative
error of order $\kOff\dt$.  This error is negligible for the manuscript
timestep and is independently testable by timestep refinement.

At finite occupancy, $\tauB$ is not constant because occupied neighboring
receptors block cross-rebinding.  The simulations therefore impose
$\phi\in(0,1)$, estimate $\tauB(\phi)$, and identify the self-consistent
ensemble occupancy as $\theta=\phi$.  Rearrangement of
Eq.~\ref{eq:renewal-balance} gives the parametric concentration-response curve
used in Figures~3--5:
\begin{equation}
 \boxed{
 C(\phi)=
 \frac{\phi}{(1-\phi)\kOn\tauB(\phi)}
 }
 \qquad (0<\phi<1).
 \label{eq:concentration-from-phi}
\end{equation}
The curve is therefore not an extrapolation of a single dilute-limit lifetime;
each point uses a retention simulation performed at the corresponding
background occupancy.

\subsection{Operational half-occupancy concentration and normalization}

The concentration producing a chosen occupancy $p$ is
\begin{equation}
 C_p=\frac{p}{(1-p)\kOn\tauB(p)}.
 \label{eq:Cp}
\end{equation}
In particular,
\begin{equation}
 \boxed{
 C_{50}=\frac{1}{\kOn\tauB(0.5)}
 },
 \qquad
 \boxed{
 \frac{C_{50}}{K_D}=\frac{1}{\kOff\tauB(0.5)}
 }.
 \label{eq:C50}
\end{equation}
The dimensionless residence amplification
\begin{equation}
 A_B(\phi)=\kOff\tauB(\phi)
 \label{eq:residence-amplification}
\end{equation}
equals unity in the classical, no-rebinding limit.  More generally,
normalizing Eq.~\ref{eq:concentration-from-phi} by
$K_D=\kOff/\kOn$ gives
\begin{equation}
 \frac{C(\phi)}{K_D}
 =\frac{\phi}{(1-\phi)\kOff\tauB(\phi)}
 =\frac{\phi}{(1-\phi)A_B(\phi)}.
 \label{eq:normalized-concentration}
\end{equation}
At the same occupancy, the concentration required by the classical model is
$C_{p,\mathrm{classical}}=K_Dp/(1-p)$.  Hence the kinetic concentration
reduction is
\begin{equation}
 \frac{C_{p,\mathrm{classical}}}{C_p}
 =\kOff\tauB(p)=A_B(p).
 \label{eq:concentration-reduction}
\end{equation}
This identity makes clear why occupancy-dependent suppression of
cross-rebinding can change the response-curve shape: if $A_B(p)$ is larger at
low $p$ than near saturation, the dilute end is shifted left more strongly
than the saturated end.

\subsection{Estimating concentration from occupancy}

If the measured fractional occupancy $\theta$ and the appropriate
occupancy-conditioned lifetime are known, the model's concentration estimate
is the inverse of Eq.~\ref{eq:effective-occupancy}:
\begin{equation}
 \boxed{
 \widehat C
 =\frac{\theta}{(1-\theta)\kOn\tauB(\theta)}
 }.
 \label{eq:estimated-concentration}
\end{equation}
This estimator is model-based: it requires the same intrinsic $\kOn$,
geometry, receptor density, and occupancy convention used to calibrate
$\tauB(\theta)$.  Near $\theta=0$ or $1$, uncertainty is amplified by the
logit factor $\theta/(1-\theta)$, so endpoint estimates require correspondingly
larger trajectory ensembles.

For receptor-resolved maps, the same equations are applied to
$\tau_{B,r}(\phi)$, the mean cumulative bound time of trajectories initiated
at receptor $r$.  This produces local operational quantities such as
\begin{equation}
 \frac{C_{50,r}}{K_D}
 =\frac{1}{\kOff\tau_{B,r}(0.5)},
 \label{eq:local-C50}
\end{equation}
which report spatially heterogeneous retention on a single structured
surface.  They are local kinetic response metrics, not receptor-specific
binding free energies.

\subsection{Censoring and statistical interpretation}

For each receptor location and condition, $\tauB$ is estimated by averaging
$T_B$ over independent trajectories and then combining independent simulation
replicates.  If a trajectory reaches $t_{\max}$ before escape, its recorded
$T_B$ is a lower bound on the eventual cumulative bound time.  The simple
sample mean used in the concentration conversion is therefore a restricted
mean and is biased downward if censoring is appreciable; this would bias
$C_{50}$ upward.  Manuscript production runs set the horizon far beyond the
intrinsic mean bound time (and kinetic sweeps scale it with $1/\kOff$), while
the censoring fraction is retained as a quality-control observable.  The
Kaplan--Meier estimator in Eq.~\ref{eq:kaplan-meier} is used for the separate
first-passage survival plots.

Equations~\ref{eq:renewal-balance}--\ref{eq:estimated-concentration} also assume
that bulk concentration is spatially equilibrated and that arrival to vacant
receptors follows the pseudo-first-order rate $\kOn C$.  The single-ligand
escape assay resolves post-dissociation retention; it does not independently
calculate macroscopic mass-transfer limitations from a finite reservoir.
